\documentclass[10pt,a4paper]{article}
\usepackage[top=1.8cm,bottom=1.8cm,left=1.9cm,right=1.9cm]{geometry}
\usepackage[T1]{fontenc}
\usepackage{lmodern}
\usepackage{microtype}
\usepackage{booktabs}
\usepackage{graphicx}
\usepackage{amsmath}
\usepackage{siunitx}
\usepackage{xcolor}
\usepackage{caption}
\usepackage{subcaption}
\usepackage{enumitem}
\usepackage[hidelinks]{hyperref}
\usepackage[numbers,sort&compress]{natbib}
\captionsetup{font=footnotesize,labelfont=bf,skip=3pt}
\setlist{nosep,leftmargin=*}
\setlength{\parskip}{2pt}
\graphicspath{{figures/}}
% generated by scripts/make_report_assets.py -- do not edit
\providecommand{\accEmbGloveMean}{}\renewcommand{\accEmbGloveMean}{87.9}
\providecommand{\accEmbRandLearnedMean}{}\renewcommand{\accEmbRandLearnedMean}{75.2}
\providecommand{\accEmbSkipgramMean}{}\renewcommand{\accEmbSkipgramMean}{83.4}
\providecommand{\accLstmGloveFrozenMean}{}\renewcommand{\accLstmGloveFrozenMean}{82.4}
\providecommand{\accLstmRandLast}{}\renewcommand{\accLstmRandLast}{62.7}
\providecommand{\accLstmRandMean}{}\renewcommand{\accLstmRandMean}{68.1}
\providecommand{\accLstmSkipgramFrozenMean}{}\renewcommand{\accLstmSkipgramFrozenMean}{79.4}
\providecommand{\accRnnRandMean}{}\renewcommand{\accRnnRandMean}{73.6}
\providecommand{\accTfidf}{}\renewcommand{\accTfidf}{89.9}
\providecommand{\accTransformerOnelRandMean}{}\renewcommand{\accTransformerOnelRandMean}{79.0}
\providecommand{\accTransformerRandLast}{}\renewcommand{\accTransformerRandLast}{73.6}
\providecommand{\accTransformerRandMean}{}\renewcommand{\accTransformerRandMean}{77.7}
\providecommand{\agInVocab}{}\renewcommand{\agInVocab}{92.5}
\providecommand{\agMeanLen}{}\renewcommand{\agMeanLen}{45}
\providecommand{\agTruncated}{}\renewcommand{\agTruncated}{0.7}
\providecommand{\agUnkBusiness}{}\renewcommand{\agUnkBusiness}{7.7}
\providecommand{\agUnkSciTech}{}\renewcommand{\agUnkSciTech}{8.6}
\providecommand{\agUnkSports}{}\renewcommand{\agUnkSports}{7.2}
\providecommand{\agUnkWorld}{}\renewcommand{\agUnkWorld}{6.5}
\providecommand{\bestEpochLongLSTM}{}\renewcommand{\bestEpochLongLSTM}{18}
\providecommand{\bestEpochLongTransformer}{}\renewcommand{\bestEpochLongTransformer}{24}
\providecommand{\bestLMRep}{}\renewcommand{\bestLMRep}{LSTM + frozen GloVe, mean}
\providecommand{\bestLMRepAcc}{}\renewcommand{\bestLMRepAcc}{82.4}
\providecommand{\bestLMRepFOne}{}\renewcommand{\bestLMRepFOne}{82.3}
\providecommand{\bestRep}{}\renewcommand{\bestRep}{TF-IDF uni+bigrams (reference)}
\providecommand{\bestRepFOne}{}\renewcommand{\bestRepFOne}{89.9}
\providecommand{\cGloveFrozenVsTrain}{}\renewcommand{\cGloveFrozenVsTrain}{7.3}
\providecommand{\cRandFrozenVsTrain}{}\renewcommand{\cRandFrozenVsTrain}{64.1}
\providecommand{\cRandVsGloveFrozen}{}\renewcommand{\cRandVsGloveFrozen}{10.2}
\providecommand{\cRandVsGloveTrain}{}\renewcommand{\cRandVsGloveTrain}{17.5}
\providecommand{\cRandVsSkipFrozen}{}\renewcommand{\cRandVsSkipFrozen}{6.1}
\providecommand{\cSkipFrozenVsTrain}{}\renewcommand{\cSkipFrozenVsTrain}{8.8}
\providecommand{\cSkipVsGloveFrozen}{}\renewcommand{\cSkipVsGloveFrozen}{4.1}
\providecommand{\cSkipVsGloveTrain}{}\renewcommand{\cSkipVsGloveTrain}{2.6}
\providecommand{\ceEosLSTM}{}\renewcommand{\ceEosLSTM}{1.30}
\providecommand{\ceEosNgram}{}\renewcommand{\ceEosNgram}{1.82}
\providecommand{\ceEosRNN}{}\renewcommand{\ceEosRNN}{1.45}
\providecommand{\ceEosTransformer}{}\renewcommand{\ceEosTransformer}{1.09}
\providecommand{\ceEosTransformerOne}{}\renewcommand{\ceEosTransformerOne}{1.27}
\providecommand{\ceMidLSTM}{}\renewcommand{\ceMidLSTM}{9.15}
\providecommand{\ceMidNgram}{}\renewcommand{\ceMidNgram}{9.77}
\providecommand{\ceMidRNN}{}\renewcommand{\ceMidRNN}{9.38}
\providecommand{\ceMidTransformer}{}\renewcommand{\ceMidTransformer}{9.05}
\providecommand{\ceMidTransformerOne}{}\renewcommand{\ceMidTransformerOne}{9.30}
\providecommand{\ceRareLSTM}{}\renewcommand{\ceRareLSTM}{13.16}
\providecommand{\ceRareNgram}{}\renewcommand{\ceRareNgram}{12.25}
\providecommand{\ceRareRNN}{}\renewcommand{\ceRareRNN}{12.89}
\providecommand{\ceRareTransformer}{}\renewcommand{\ceRareTransformer}{12.74}
\providecommand{\ceRareTransformerOne}{}\renewcommand{\ceRareTransformerOne}{13.05}
\providecommand{\ceTopLSTM}{}\renewcommand{\ceTopLSTM}{2.26}
\providecommand{\ceTopNgram}{}\renewcommand{\ceTopNgram}{2.65}
\providecommand{\ceTopRNN}{}\renewcommand{\ceTopRNN}{2.48}
\providecommand{\ceTopTransformer}{}\renewcommand{\ceTopTransformer}{2.40}
\providecommand{\ceTopTransformerOne}{}\renewcommand{\ceTopTransformerOne}{2.59}
\providecommand{\ceUnkLSTM}{}\renewcommand{\ceUnkLSTM}{3.19}
\providecommand{\ceUnkNgram}{}\renewcommand{\ceUnkNgram}{3.35}
\providecommand{\ceUnkRNN}{}\renewcommand{\ceUnkRNN}{3.35}
\providecommand{\ceUnkTransformer}{}\renewcommand{\ceUnkTransformer}{3.51}
\providecommand{\ceUnkTransformerOne}{}\renewcommand{\ceUnkTransformerOne}{3.63}
\providecommand{\confBusinessSciTech}{}\renewcommand{\confBusinessSciTech}{526}
\providecommand{\confWorldBusiness}{}\renewcommand{\confWorldBusiness}{285}
\providecommand{\confWorldSports}{}\renewcommand{\confWorldSports}{170}
\providecommand{\confusionRep}{}\renewcommand{\confusionRep}{LSTM + frozen GloVe, mean}
\providecommand{\coreMLSTM}{}\renewcommand{\coreMLSTM}{1.44}
\providecommand{\coreMNgram}{}\renewcommand{\coreMNgram}{--}
\providecommand{\coreMRNN}{}\renewcommand{\coreMRNN}{0.36}
\providecommand{\coreMTransformer}{}\renewcommand{\coreMTransformer}{4.37}
\providecommand{\coreMTransformerOne}{}\renewcommand{\coreMTransformerOne}{1.12}
\providecommand{\deltaBlocksLSTM}{}\renewcommand{\deltaBlocksLSTM}{16.2}
\providecommand{\deltaBlocksRNN}{}\renewcommand{\deltaBlocksRNN}{4.3}
\providecommand{\deltaBlocksTransformer}{}\renewcommand{\deltaBlocksTransformer}{5.6}
\providecommand{\deltaOneLayerVsLSTM}{}\renewcommand{\deltaOneLayerVsLSTM}{37.2}
\providecommand{\deltaTransformerVsLSTM}{}\renewcommand{\deltaTransformerVsLSTM}{6.6}
\providecommand{\distinctTopkLSTM}{}\renewcommand{\distinctTopkLSTM}{0.95}
\providecommand{\distinctTopkLSTMGloveFrozen}{}\renewcommand{\distinctTopkLSTMGloveFrozen}{0.92}
\providecommand{\distinctTopkLSTMGloveTrain}{}\renewcommand{\distinctTopkLSTMGloveTrain}{0.95}
\providecommand{\distinctTopkLSTMRand}{}\renewcommand{\distinctTopkLSTMRand}{0.95}
\providecommand{\distinctTopkLSTMRandFrozen}{}\renewcommand{\distinctTopkLSTMRandFrozen}{0.84}
\providecommand{\distinctTopkLSTMSkipFrozen}{}\renewcommand{\distinctTopkLSTMSkipFrozen}{0.91}
\providecommand{\distinctTopkLSTMSkipTrain}{}\renewcommand{\distinctTopkLSTMSkipTrain}{0.90}
\providecommand{\distinctTopkNgram}{}\renewcommand{\distinctTopkNgram}{0.98}
\providecommand{\distinctTopkRNN}{}\renewcommand{\distinctTopkRNN}{0.91}
\providecommand{\distinctTopkTransformer}{}\renewcommand{\distinctTopkTransformer}{0.94}
\providecommand{\distinctTopkTransformerOne}{}\renewcommand{\distinctTopkTransformerOne}{0.87}
\providecommand{\epochReachLSTMGloveFrozen}{}\renewcommand{\epochReachLSTMGloveFrozen}{8}
\providecommand{\epochReachLSTMGloveTrain}{}\renewcommand{\epochReachLSTMGloveTrain}{7}
\providecommand{\epochReachLSTMRandFrozen}{}\renewcommand{\epochReachLSTMRandFrozen}{never}
\providecommand{\epochReachLSTMSkipFrozen}{}\renewcommand{\epochReachLSTMSkipFrozen}{10}
\providecommand{\epochReachLSTMSkipTrain}{}\renewcommand{\epochReachLSTMSkipTrain}{8}
\providecommand{\epochSecLSTM}{}\renewcommand{\epochSecLSTM}{110}
\providecommand{\epochSecRNN}{}\renewcommand{\epochSecRNN}{116}
\providecommand{\epochSecTransformer}{}\renewcommand{\epochSecTransformer}{149}
\providecommand{\epochSecTransformerOne}{}\renewcommand{\epochSecTransformerOne}{110}
\providecommand{\epochsRunLongLSTM}{}\renewcommand{\epochsRunLongLSTM}{21}
\providecommand{\epochsRunLongTransformer}{}\renewcommand{\epochsRunLongTransformer}{24}
\providecommand{\fOneEmbGloveMean}{}\renewcommand{\fOneEmbGloveMean}{87.9}
\providecommand{\fOneEmbRandLearnedMean}{}\renewcommand{\fOneEmbRandLearnedMean}{75.1}
\providecommand{\fOneEmbSkipgramMean}{}\renewcommand{\fOneEmbSkipgramMean}{83.4}
\providecommand{\fOneLstmGloveFrozenMean}{}\renewcommand{\fOneLstmGloveFrozenMean}{82.3}
\providecommand{\fOneLstmRandLast}{}\renewcommand{\fOneLstmRandLast}{62.3}
\providecommand{\fOneLstmRandMean}{}\renewcommand{\fOneLstmRandMean}{67.9}
\providecommand{\fOneLstmSkipgramFrozenMean}{}\renewcommand{\fOneLstmSkipgramFrozenMean}{79.3}
\providecommand{\fOneRnnRandMean}{}\renewcommand{\fOneRnnRandMean}{73.5}
\providecommand{\fOneSmallEmbGloveMean}{}\renewcommand{\fOneSmallEmbGloveMean}{85.6}
\providecommand{\fOneSmallEmbRandLearnedMean}{}\renewcommand{\fOneSmallEmbRandLearnedMean}{63.3}
\providecommand{\fOneSmallEmbSkipgramMean}{}\renewcommand{\fOneSmallEmbSkipgramMean}{80.1}
\providecommand{\fOneSmallLstmGloveFrozenMean}{}\renewcommand{\fOneSmallLstmGloveFrozenMean}{77.8}
\providecommand{\fOneSmallLstmRandLast}{}\renewcommand{\fOneSmallLstmRandLast}{54.5}
\providecommand{\fOneSmallLstmRandMean}{}\renewcommand{\fOneSmallLstmRandMean}{60.6}
\providecommand{\fOneSmallLstmSkipgramFrozenMean}{}\renewcommand{\fOneSmallLstmSkipgramFrozenMean}{75.2}
\providecommand{\fOneSmallRnnRandMean}{}\renewcommand{\fOneSmallRnnRandMean}{67.2}
\providecommand{\fOneSmallTfidf}{}\renewcommand{\fOneSmallTfidf}{78.3}
\providecommand{\fOneSmallTransformerOnelRandMean}{}\renewcommand{\fOneSmallTransformerOnelRandMean}{70.9}
\providecommand{\fOneSmallTransformerRandLast}{}\renewcommand{\fOneSmallTransformerRandLast}{66.6}
\providecommand{\fOneSmallTransformerRandMean}{}\renewcommand{\fOneSmallTransformerRandMean}{69.5}
\providecommand{\fOneTfidf}{}\renewcommand{\fOneTfidf}{89.9}
\providecommand{\fOneTransformerOnelRandMean}{}\renewcommand{\fOneTransformerOnelRandMean}{78.9}
\providecommand{\fOneTransformerRandLast}{}\renewcommand{\fOneTransformerRandLast}{73.6}
\providecommand{\fOneTransformerRandMean}{}\renewcommand{\fOneTransformerRandMean}{77.6}
\providecommand{\gainMidLSTMGloveFrozen}{}\renewcommand{\gainMidLSTMGloveFrozen}{0.22}
\providecommand{\gainMidLSTMGloveTrain}{}\renewcommand{\gainMidLSTMGloveTrain}{0.26}
\providecommand{\gainMidLSTMRand}{}\renewcommand{\gainMidLSTMRand}{0.00}
\providecommand{\gainMidLSTMRandFrozen}{}\renewcommand{\gainMidLSTMRandFrozen}{-0.58}
\providecommand{\gainMidLSTMSkipFrozen}{}\renewcommand{\gainMidLSTMSkipFrozen}{0.19}
\providecommand{\gainMidLSTMSkipTrain}{}\renewcommand{\gainMidLSTMSkipTrain}{0.25}
\providecommand{\gainRareLSTMGloveFrozen}{}\renewcommand{\gainRareLSTMGloveFrozen}{0.47}
\providecommand{\gainRareLSTMGloveTrain}{}\renewcommand{\gainRareLSTMGloveTrain}{0.27}
\providecommand{\gainRareLSTMRand}{}\renewcommand{\gainRareLSTMRand}{0.00}
\providecommand{\gainRareLSTMRandFrozen}{}\renewcommand{\gainRareLSTMRandFrozen}{-0.03}
\providecommand{\gainRareLSTMSkipFrozen}{}\renewcommand{\gainRareLSTMSkipFrozen}{0.50}
\providecommand{\gainRareLSTMSkipTrain}{}\renewcommand{\gainRareLSTMSkipTrain}{0.42}
\providecommand{\gainTopLSTMGloveFrozen}{}\renewcommand{\gainTopLSTMGloveFrozen}{0.00}
\providecommand{\gainTopLSTMGloveTrain}{}\renewcommand{\gainTopLSTMGloveTrain}{+0.08}
\providecommand{\gainTopLSTMRand}{}\renewcommand{\gainTopLSTMRand}{0.00}
\providecommand{\gainTopLSTMRandFrozen}{}\renewcommand{\gainTopLSTMRandFrozen}{-0.21}
\providecommand{\gainTopLSTMSkipFrozen}{}\renewcommand{\gainTopLSTMSkipFrozen}{-0.05}
\providecommand{\gainTopLSTMSkipTrain}{}\renewcommand{\gainTopLSTMSkipTrain}{+0.05}
\providecommand{\gapLSTMGloveFrozen}{}\renewcommand{\gapLSTMGloveFrozen}{13.7}
\providecommand{\gapLSTMGloveTrain}{}\renewcommand{\gapLSTMGloveTrain}{26.4}
\providecommand{\gapLSTMRand}{}\renewcommand{\gapLSTMRand}{32.2}
\providecommand{\gapLSTMRandFrozen}{}\renewcommand{\gapLSTMRandFrozen}{18.5}
\providecommand{\gapLSTMSkipFrozen}{}\renewcommand{\gapLSTMSkipFrozen}{14.7}
\providecommand{\gapLSTMSkipTrain}{}\renewcommand{\gapLSTMSkipTrain}{29.8}
\providecommand{\genMsLSTM}{}\renewcommand{\genMsLSTM}{1.7}
\providecommand{\genMsNgram}{}\renewcommand{\genMsNgram}{19.7}
\providecommand{\genMsRNN}{}\renewcommand{\genMsRNN}{1.6}
\providecommand{\genMsTransformer}{}\renewcommand{\genMsTransformer}{3.4}
\providecommand{\genMsTransformerOne}{}\renewcommand{\genMsTransformerOne}{2.3}
\providecommand{\gloveCoverageTokens}{}\renewcommand{\gloveCoverageTokens}{94.74}
\providecommand{\gloveCoverageTypes}{}\renewcommand{\gloveCoverageTypes}{96.2}
\providecommand{\gloveMissing}{}\renewcommand{\gloveMissing}{1,104}
\providecommand{\gloveMissingExamples}{}\renewcommand{\gloveMissingExamples}{\textless{}unk\textgreater{}, \textless{}eos\textgreater{}, @-@, @.@, @,@, 't}
\providecommand{\gnMaxLSTMGloveFrozen}{}\renewcommand{\gnMaxLSTMGloveFrozen}{1.3}
\providecommand{\gnMaxLSTMGloveFrozenSeedOne}{}\renewcommand{\gnMaxLSTMGloveFrozenSeedOne}{1.4}
\providecommand{\gnMaxLSTMGloveTrain}{}\renewcommand{\gnMaxLSTMGloveTrain}{2.3}
\providecommand{\gnMaxLSTMRand}{}\renewcommand{\gnMaxLSTMRand}{1.8}
\providecommand{\gnMaxLSTMRandFrozen}{}\renewcommand{\gnMaxLSTMRandFrozen}{5.8}
\providecommand{\gnMaxLSTMRandSeedOne}{}\renewcommand{\gnMaxLSTMRandSeedOne}{13.1}
\providecommand{\gnMaxLSTMSkipFrozen}{}\renewcommand{\gnMaxLSTMSkipFrozen}{1.3}
\providecommand{\gnMaxLSTMSkipFrozenSeedOne}{}\renewcommand{\gnMaxLSTMSkipFrozenSeedOne}{1.3}
\providecommand{\gnMaxLSTMSkipTrain}{}\renewcommand{\gnMaxLSTMSkipTrain}{1.3}
\providecommand{\gnMeanLSTMGloveFrozen}{}\renewcommand{\gnMeanLSTMGloveFrozen}{0.28}
\providecommand{\gnMeanLSTMGloveTrain}{}\renewcommand{\gnMeanLSTMGloveTrain}{0.27}
\providecommand{\gnMeanLSTMRand}{}\renewcommand{\gnMeanLSTMRand}{0.27}
\providecommand{\gnMeanLSTMRandFrozen}{}\renewcommand{\gnMeanLSTMRandFrozen}{0.32}
\providecommand{\gnMeanLSTMSkipFrozen}{}\renewcommand{\gnMeanLSTMSkipFrozen}{0.29}
\providecommand{\gnMeanLSTMSkipTrain}{}\renewcommand{\gnMeanLSTMSkipTrain}{0.27}
\providecommand{\nLabelsFull}{}\renewcommand{\nLabelsFull}{20,000}
\providecommand{\nLabelsSmall}{}\renewcommand{\nLabelsSmall}{500}
\providecommand{\nTestDocs}{}\renewcommand{\nTestDocs}{7,600}
\providecommand{\nTestTargets}{}\renewcommand{\nTestTargets}{244,102}
\providecommand{\nTrainTok}{}\renewcommand{\nTrainTok}{2,075,678}
\providecommand{\nValTok}{}\renewcommand{\nValTok}{216,348}
\providecommand{\nVocab}{}\renewcommand{\nVocab}{28,913}
\providecommand{\pplBlocksLSTM}{}\renewcommand{\pplBlocksLSTM}{150.4}
\providecommand{\pplBlocksNgram}{}\renewcommand{\pplBlocksNgram}{--}
\providecommand{\pplBlocksRNN}{}\renewcommand{\pplBlocksRNN}{177.2}
\providecommand{\pplBlocksTransformer}{}\renewcommand{\pplBlocksTransformer}{146.4}
\providecommand{\pplBlocksTransformerOne}{}\renewcommand{\pplBlocksTransformerOne}{174.7}
\providecommand{\pplLSTM}{}\renewcommand{\pplLSTM}{134.2}
\providecommand{\pplLSTMGloveFrozen}{}\renewcommand{\pplLSTMGloveFrozen}{123.8}
\providecommand{\pplLSTMGloveTrain}{}\renewcommand{\pplLSTMGloveTrain}{116.5}
\providecommand{\pplLSTMRand}{}\renewcommand{\pplLSTMRand}{134.0}
\providecommand{\pplLSTMRandFrozen}{}\renewcommand{\pplLSTMRandFrozen}{198.0}
\providecommand{\pplLSTMSkipFrozen}{}\renewcommand{\pplLSTMSkipFrozen}{127.9}
\providecommand{\pplLSTMSkipTrain}{}\renewcommand{\pplLSTMSkipTrain}{119.0}
\providecommand{\pplLongBlocksLSTM}{}\renewcommand{\pplLongBlocksLSTM}{141.0}
\providecommand{\pplLongBlocksTransformer}{}\renewcommand{\pplLongBlocksTransformer}{137.3}
\providecommand{\pplLongLSTM}{}\renewcommand{\pplLongLSTM}{124.6}
\providecommand{\pplLongNgram}{}\renewcommand{\pplLongNgram}{--}
\providecommand{\pplLongRNN}{}\renewcommand{\pplLongRNN}{--}
\providecommand{\pplLongTransformer}{}\renewcommand{\pplLongTransformer}{132.3}
\providecommand{\pplLongTransformerOne}{}\renewcommand{\pplLongTransformerOne}{--}
\providecommand{\pplNgram}{}\renewcommand{\pplNgram}{214.9}
\providecommand{\pplRNN}{}\renewcommand{\pplRNN}{173.0}
\providecommand{\pplRangeLSTMGloveFrozen}{}\renewcommand{\pplRangeLSTMGloveFrozen}{0.4}
\providecommand{\pplRangeLSTMRand}{}\renewcommand{\pplRangeLSTMRand}{0.3}
\providecommand{\pplRangeLSTMSkipFrozen}{}\renewcommand{\pplRangeLSTMSkipFrozen}{0.5}
\providecommand{\pplSeedOneLSTMGloveFrozen}{}\renewcommand{\pplSeedOneLSTMGloveFrozen}{124.2}
\providecommand{\pplSeedOneLSTMRand}{}\renewcommand{\pplSeedOneLSTMRand}{133.7}
\providecommand{\pplSeedOneLSTMSkipFrozen}{}\renewcommand{\pplSeedOneLSTMSkipFrozen}{127.3}
\providecommand{\pplTransformer}{}\renewcommand{\pplTransformer}{140.9}
\providecommand{\pplTransformerOne}{}\renewcommand{\pplTransformerOne}{171.4}
\providecommand{\repGreedyLSTM}{}\renewcommand{\repGreedyLSTM}{25}
\providecommand{\repGreedyLSTMGloveFrozen}{}\renewcommand{\repGreedyLSTMGloveFrozen}{46}
\providecommand{\repGreedyLSTMGloveTrain}{}\renewcommand{\repGreedyLSTMGloveTrain}{46}
\providecommand{\repGreedyLSTMRand}{}\renewcommand{\repGreedyLSTMRand}{25}
\providecommand{\repGreedyLSTMRandFrozen}{}\renewcommand{\repGreedyLSTMRandFrozen}{56}
\providecommand{\repGreedyLSTMSkipFrozen}{}\renewcommand{\repGreedyLSTMSkipFrozen}{41}
\providecommand{\repGreedyLSTMSkipTrain}{}\renewcommand{\repGreedyLSTMSkipTrain}{41}
\providecommand{\repGreedyNgram}{}\renewcommand{\repGreedyNgram}{58}
\providecommand{\repGreedyRNN}{}\renewcommand{\repGreedyRNN}{74}
\providecommand{\repGreedyTransformer}{}\renewcommand{\repGreedyTransformer}{31}
\providecommand{\repGreedyTransformerOne}{}\renewcommand{\repGreedyTransformerOne}{57}
\providecommand{\repTopkLSTM}{}\renewcommand{\repTopkLSTM}{0.0}
\providecommand{\repTopkNgram}{}\renewcommand{\repTopkNgram}{0.0}
\providecommand{\repTopkRNN}{}\renewcommand{\repTopkRNN}{1.3}
\providecommand{\repTopkTransformer}{}\renewcommand{\repTopkTransformer}{0.0}
\providecommand{\repTopkTransformerOne}{}\renewcommand{\repTopkTransformerOne}{2.6}
\providecommand{\rowNormGloveTable}{}\renewcommand{\rowNormGloveTable}{6.47}
\providecommand{\rowNormSkipgramTable}{}\renewcommand{\rowNormSkipgramTable}{2.37}
\providecommand{\scoreTokSLSTM}{}\renewcommand{\scoreTokSLSTM}{26,333}
\providecommand{\scoreTokSNgram}{}\renewcommand{\scoreTokSNgram}{318,385}
\providecommand{\scoreTokSRNN}{}\renewcommand{\scoreTokSRNN}{6,779}
\providecommand{\scoreTokSTransformer}{}\renewcommand{\scoreTokSTransformer}{12,694}
\providecommand{\scoreTokSTransformerOne}{}\renewcommand{\scoreTokSTransformerOne}{18,758}
\providecommand{\skipgramMinutes}{}\renewcommand{\skipgramMinutes}{1.9}
\providecommand{\spearmanN}{}\renewcommand{\spearmanN}{6}
\providecommand{\spearmanRho}{}\renewcommand{\spearmanRho}{-0.60}
\providecommand{\trainMinLSTM}{}\renewcommand{\trainMinLSTM}{22.1}
\providecommand{\trainMinLongLSTM}{}\renewcommand{\trainMinLongLSTM}{41}
\providecommand{\trainMinLongTransformer}{}\renewcommand{\trainMinLongTransformer}{62}
\providecommand{\trainMinNgram}{}\renewcommand{\trainMinNgram}{0.1}
\providecommand{\trainMinRNN}{}\renewcommand{\trainMinRNN}{23.2}
\providecommand{\trainMinTransformer}{}\renewcommand{\trainMinTransformer}{29.8}
\providecommand{\trainMinTransformerOne}{}\renewcommand{\trainMinTransformerOne}{22.0}
\providecommand{\trainPplLastLSTM}{}\renewcommand{\trainPplLastLSTM}{113.1}
\providecommand{\trainPplLastRNN}{}\renewcommand{\trainPplLastRNN}{144.9}
\providecommand{\trainPplLastTransformer}{}\renewcommand{\trainPplLastTransformer}{135.0}
\providecommand{\trainPplLastTransformerOne}{}\renewcommand{\trainPplLastTransformerOne}{168.1}
\providecommand{\unkFracTest}{}\renewcommand{\unkFracTest}{6.2}
\providecommand{\unkRateCorrect}{}\renewcommand{\unkRateCorrect}{7.5}
\providecommand{\unkRateWrong}{}\renewcommand{\unkRateWrong}{7.4}
\providecommand{\valBudgetLSTM}{}\renewcommand{\valBudgetLSTM}{214.8}
\providecommand{\valBudgetNgram}{}\renewcommand{\valBudgetNgram}{--}
\providecommand{\valBudgetRNN}{}\renewcommand{\valBudgetRNN}{210.0}
\providecommand{\valBudgetTransformer}{}\renewcommand{\valBudgetTransformer}{227.9}
\providecommand{\valBudgetTransformerOne}{}\renewcommand{\valBudgetTransformerOne}{234.0}
\providecommand{\valEpThreeLSTM}{}\renewcommand{\valEpThreeLSTM}{237.3}
\providecommand{\valEpThreeLSTMGloveFrozen}{}\renewcommand{\valEpThreeLSTMGloveFrozen}{205.7}
\providecommand{\valEpThreeLSTMGloveTrain}{}\renewcommand{\valEpThreeLSTMGloveTrain}{190.1}
\providecommand{\valEpThreeLSTMRand}{}\renewcommand{\valEpThreeLSTMRand}{237.3}
\providecommand{\valEpThreeLSTMRandFrozen}{}\renewcommand{\valEpThreeLSTMRandFrozen}{358.1}
\providecommand{\valEpThreeLSTMSkipFrozen}{}\renewcommand{\valEpThreeLSTMSkipFrozen}{228.6}
\providecommand{\valEpThreeLSTMSkipTrain}{}\renewcommand{\valEpThreeLSTMSkipTrain}{200.8}
\providecommand{\valEpThreeNgram}{}\renewcommand{\valEpThreeNgram}{--}
\providecommand{\valEpThreeRNN}{}\renewcommand{\valEpThreeRNN}{219.6}
\providecommand{\valEpThreeTransformer}{}\renewcommand{\valEpThreeTransformer}{227.9}
\providecommand{\valEpThreeTransformerOne}{}\renewcommand{\valEpThreeTransformerOne}{256.9}
\providecommand{\valEpTwelveLSTMGloveFrozen}{}\renewcommand{\valEpTwelveLSTMGloveFrozen}{131.8}
\providecommand{\valEpTwelveLSTMGloveFrozenSeedOne}{}\renewcommand{\valEpTwelveLSTMGloveFrozenSeedOne}{132.3}
\providecommand{\valEpTwelveLSTMGloveTrain}{}\renewcommand{\valEpTwelveLSTMGloveTrain}{125.0}
\providecommand{\valEpTwelveLSTMRand}{}\renewcommand{\valEpTwelveLSTMRand}{145.4}
\providecommand{\valEpTwelveLSTMRandFrozen}{}\renewcommand{\valEpTwelveLSTMRandFrozen}{218.8}
\providecommand{\valEpTwelveLSTMRandSeedOne}{}\renewcommand{\valEpTwelveLSTMRandSeedOne}{143.9}
\providecommand{\valEpTwelveLSTMSkipFrozen}{}\renewcommand{\valEpTwelveLSTMSkipFrozen}{136.6}
\providecommand{\valEpTwelveLSTMSkipFrozenSeedOne}{}\renewcommand{\valEpTwelveLSTMSkipFrozenSeedOne}{136.2}
\providecommand{\valEpTwelveLSTMSkipTrain}{}\renewcommand{\valEpTwelveLSTMSkipTrain}{126.9}
\providecommand{\valPplLastLSTM}{}\renewcommand{\valPplLastLSTM}{145.4}
\providecommand{\valPplLastRNN}{}\renewcommand{\valPplLastRNN}{187.3}
\providecommand{\valPplLastTransformer}{}\renewcommand{\valPplLastTransformer}{150.1}
\providecommand{\valPplLastTransformerOne}{}\renewcommand{\valPplLastTransformerOne}{183.1}
\providecommand{\appendix}{\textit{[pending]}}
\providecommand{\approx}{\textit{[pending]}}
\providecommand{\author}{\textit{[pending]}}
\providecommand{\begin}{\textit{[pending]}}
\providecommand{\bibliography}{\textit{[pending]}}
\providecommand{\bibliographystyle}{\textit{[pending]}}
\providecommand{\caption}{\textit{[pending]}}
\providecommand{\captionsetup}{\textit{[pending]}}
\providecommand{\centering}{\textit{[pending]}}
\providecommand{\citep}{\textit{[pending]}}
\providecommand{\clearpage}{\textit{[pending]}}
\providecommand{\csname}{\textit{[pending]}}
\providecommand{\date}{\textit{[pending]}}
\providecommand{\detokenize}{\textit{[pending]}}
\providecommand{\documentclass}{\textit{[pending]}}
\providecommand{\emph}{\textit{[pending]}}
\providecommand{\endcsname}{\textit{[pending]}}
\providecommand{\expandafter}{\textit{[pending]}}
\providecommand{\fbox}{\textit{[pending]}}
\providecommand{\graphicspath}{\textit{[pending]}}
\providecommand{\hfill}{\textit{[pending]}}
\providecommand{\includegraphics}{\textit{[pending]}}
\providecommand{\input}{\textit{[pending]}}
\providecommand{\label}{\textit{[pending]}}
\providecommand{\large}{\textit{[pending]}}
\providecommand{\linewidth}{\textit{[pending]}}
\providecommand{\makeatletter}{\textit{[pending]}}
\providecommand{\makeatother}{\textit{[pending]}}
\providecommand{\maketitle}{\textit{[pending]}}
\providecommand{\mathcal}{\textit{[pending]}}
\providecommand{\newcommand}{\textit{[pending]}}
\providecommand{\paragraph}{\textit{[pending]}}
\providecommand{\parskip}{\textit{[pending]}}
\providecommand{\quad}{\textit{[pending]}}
\providecommand{\scriptsize}{\textit{[pending]}}
\providecommand{\section}{\textit{[pending]}}
\providecommand{\setlength}{\textit{[pending]}}
\providecommand{\setlist}{\textit{[pending]}}
\providecommand{\sigma}{\textit{[pending]}}
\providecommand{\small}{\textit{[pending]}}
\providecommand{\textbf}{\textit{[pending]}}
\providecommand{\textgreater}{\textit{[pending]}}
\providecommand{\textit}{\textit{[pending]}}
\providecommand{\texttt}{\textit{[pending]}}
\providecommand{\times}{\textit{[pending]}}
\providecommand{\title}{\textit{[pending]}}
\providecommand{\usepackage}{\textit{[pending]}}
\providecommand{\vspace}{\textit{[pending]}}
   % macros generated by scripts/make_report_assets.py
% defaults so the document compiles before every experiment has run
\makeatletter
\newcommand{\pendingmacro}[1]{\@ifundefined{#1}{\expandafter\newcommand\csname #1\endcsname{\textit{[pending]}}}{}}
\makeatother
\pendingmacro{confusionRep}\pendingmacro{bestRep}\pendingmacro{bestRepFOne}\pendingmacro{bestLMRep}\pendingmacro{bestLMRepFOne}
\pendingmacro{spearmanRho}\pendingmacro{spearmanN}\pendingmacro{unkRateCorrect}\pendingmacro{unkRateWrong}
\pendingmacro{pplNgram}\pendingmacro{pplRNN}\pendingmacro{pplLSTM}\pendingmacro{pplTransformer}\pendingmacro{pplTransformerOne}
\pendingmacro{gloveCoverageTypes}\pendingmacro{gloveCoverageTokens}\pendingmacro{gloveMissing}\pendingmacro{skipgramMinutes}
\pendingmacro{rowNormGloveTable}\pendingmacro{rowNormSkipgramTable}\pendingmacro{gloveMissingExamples}\pendingmacro{pplLongLSTM}\pendingmacro{pplLongTransformer}

\title{\vspace{-1.2cm}DSA4213 Assignment 1: Does better next-token prediction give better representations?\\
\large Language models, embeddings and downstream topic classification on WikiText-2 and AG News}
\author{\textbf{Aniroodh Chaudhary} \quad A0284770W \quad AY2026/27 Semester 1}
\date{\vspace{-0.4cm}}

\begin{document}
\maketitle
\vspace{-0.8cm}

\section{Question and setup}
Two things are commonly assumed about small language models (LMs): that lower perplexity means a better
model, and that better LMs give better text representations. This report tests both on one corpus.
The comparison question is: \emph{does better next-token prediction translate into better frozen
representations for topic classification, and how much do input-embedding choices matter?}
Part~I trains four LM families on WikiText-2 under one recipe; Part~II keeps the LSTM and
varies only its input-embedding table in a $3\times2$ (source $\times$ trainability) design;
Part~III feeds the frozen models' states to one linear classifier on AG News. Everything not under
study is held fixed: vocabulary, token stream, embedding width (300), output layer (untied),
optimiser, schedule, early stopping, evaluation targets, prompts, decoding and classifier.

\section{Data and common protocol}
\paragraph{Corpus and vocabulary.}
WikiText-2 \citep{merity2017pointer} (pre-tokenised, with its own \texttt{<unk>}) is lowercased
once, one \texttt{<eos>} is appended per line, and the vocabulary is built from the train split
only: \nVocab{} types. Lowercasing is a deliberate fairness decision: GloVe-6B is uncased, so a
cased vocabulary would let \emph{The}/\emph{the} share one vector only in the GloVe setting of
Part~II. It shrinks the vocabulary from 33,278 types, so the perplexities below are not comparable
with published WikiText-2 numbers. Validation and test contain no out-of-vocabulary tokens
beyond the dataset's own \texttt{<unk>} (\unkFracTest\% of test tokens), and \texttt{<unk>} is an
ordinary vocabulary item for every model, including the $n$-gram.

\paragraph{Streams and perplexity.}
Each split is one continuous token stream (train \nTrainTok{} tokens, validation \nValTok);
no padding exists, so nothing has to be excluded from any loss. Every model is scored on the
same \nTestTargets{} test targets, and perplexity is $\exp$ of the mean per-target cross-entropy.
Recurrent models are scored with batch~1 and the hidden state carried over the entire stream
(full history); the Transformer with a sliding window of 128 tokens moved by 32, scoring the last
32 targets of each window (each target sees $\ge$96 tokens of context). A second protocol,
non-overlapping 128-token blocks with the state reset (identical for every model), isolates the
effect of context length. The same protocol is used for validation during training, so the
curves and the final numbers are directly comparable.

\paragraph{Models and recipe.}
All neural models share a $300$-d input embedding initialised $\mathcal{U}(-0.1,0.1)$ and a
$300\!\times\!V$ output layer, so they differ only in the sequence core: a 2-layer tanh RNN
(0.36\,M core parameters), a 2-layer LSTM (1.44\,M), a 4-layer pre-LN Transformer with 6 heads and
FFN width 1200 (4.37\,M), and, as a parameter-count control, the same Transformer with 1 layer
(1.09\,M, close to the LSTM). The $n$-gram baseline is interpolated Kneser--Ney
\citep{kneser1995improved,chen1998empirical} with one discount per order, orders 2--5 fitted and
the order chosen on validation. The neural recipe is identical for every run: batch 32 $\times$
128-token BPTT chunks with the state detached between chunks, dropout 0.3, Adam ($10^{-3}$,
200-step warm-up), learning rate halved after one epoch without validation improvement, early
stopping after three, at most 12 epochs, global gradient-norm clipping at 1.0 with the pre-clip
norm logged at every step. A unit test perturbs future tokens and checks that no model's logits at
earlier positions change (causality), and that input and output tables are not tied.

\paragraph{Cost measurement.}
All runs use one Apple M4 (10 cores, 16\,GB) through PyTorch's MPS backend, one process at a time.
Timers synchronise the device before and after each measurement; epoch~1 (shader compilation) is
reported separately from later epochs. Model size is the parameter count and the saved fp32
checkpoint. The common inference workload is (A)~scoring the full test stream with the protocol
above and (B)~generating 50 tokens for each of 5 prompts at batch~1 with top-$k$ sampling (the
Transformer without a KV cache). MPS is not bit-deterministic; the seed-1 repeats in Part~II give
the run-to-run spread, about one perplexity point, and \texttt{--device cpu} reproduces exactly.

\paragraph{Generation.}
Five prompts are the first 10 tokens of seeded-random \emph{validation} lines (never test), and
every model continues each prompt for 50 tokens with greedy decoding and with top-$k$=40, $T$=1
sampling from one shared CPU sampler, so decoding is matched exactly across the $n$-gram and
neural models. Continuations are summarised by distinct-$n$ ratios and the fraction of 4-grams
that repeat an earlier 4-gram of the same continuation.

\newcommand{\inputifexists}[1]{\IfFileExists{#1}{\input{#1}}{\textit{[\texttt{\detokenize{#1}} pending]}}}
\newcommand{\figifexists}[2]{\IfFileExists{figures/#1}{\includegraphics[#2]{#1}}{\fbox{\textit{figure \texttt{\detokenize{#1}} pending}}}}

\section{Part I: four language models on one corpus}
\begin{table}[t]\centering\small
\inputifexists{tables/part1_quality.tex}
\caption{Part I prediction quality on the common \nTestTargets{} test targets. ``ctx$\le$128'' resets the
state every 128 tokens for every model; the epoch-3 and 8-minute columns compare the models at a
matched number of updates and at a matched wall-clock budget; the last column is a separate 24-epoch run of the same configuration (convergence check).}
\label{tab:p1q}
\end{table}
\begin{table}[t]\centering\small
\inputifexists{tables/part1_cost.tex}
\caption{Part I cost on one Apple M4 (MPS). ``Core'' excludes the identical embedding and output layers.
Score: test-stream scoring throughput under the reported protocol; generate: batch-1 sampling.}
\label{tab:p1c}
\end{table}
\begin{figure}[t]\centering
\figifexists{part1_curves.pdf}{width=\linewidth}
\caption{Validation perplexity (solid) and train perplexity (dotted) per epoch (left) and validation perplexity against training wall-clock (right); dashed: the separate 24-epoch runs.}
\label{fig:p1curves}
\end{figure}
\begin{figure}[t]\centering
\figifexists{part1_buckets.pdf}{width=\linewidth}
\caption{Test cross-entropy split by how often the target word occurs in the train split.}
\label{fig:p1buckets}
\end{figure}
\begin{table}[t]\centering\scriptsize
\inputifexists{tables/part1_generations.tex}
\caption{Matched top-$k$ continuations of one validation prompt (all prompts and greedy decoding in Appendix~A).
distinct-2: unique bigrams / bigrams; rep.\ 4-gram: fraction of 4-grams repeating an earlier one.}
\label{tab:p1gen}
\end{table}
\paragraph{Observations.}
\emph{Prediction quality} (Table~\ref{tab:p1q}). With full history the order is LSTM (\pplLSTM) $<$
Transformer (\pplTransformer) $<$ RNN (\pplRNN) $\approx$ 1-layer Transformer (\pplTransformerOne) $<$
KN 5-gram (\pplNgram). The order of the two best models \emph{reverses} when every model's context is
reset every 128 tokens: the Transformer loses only \deltaBlocksTransformer{} points, the LSTM loses \deltaBlocksLSTM{} and ends up behind it.
\emph{Learning behaviour} (Fig.~\ref{fig:p1curves}). All four neural models were still improving at the
12-epoch cap: no learning-rate halving and no early stop fired, and the train/validation gap grows
(LSTM: \trainPplLastLSTM{} vs.\ \valPplLastLSTM{} at epoch 12) without validation loss turning up, i.e.\ no
overfitting at this budget. Per epoch the LSTM and 4-layer Transformer track each other closely; per
wall-clock minute the LSTM leads because its epoch takes \epochSecLSTM{}\,s against \epochSecTransformer{}\,s.
After three epochs the tanh RNN is the \emph{best} model (\valEpThreeRNN{} vs.\ \valEpThreeLSTM{} and
\valEpThreeTransformer); it learns fastest and saturates first.
\emph{Where the differences sit} (Fig.~\ref{fig:p1buckets}). On the 3\% of test targets seen at most 10
times in training the $n$-gram is the best model (\ceRareNgram{} nats vs.\ \ceRareLSTM{} for the LSTM); every
neural model wins on words seen $\ge$11 times and the margin widens with frequency (\textgreater10k bucket:
LSTM \ceTopLSTM{} vs.\ $n$-gram \ceTopNgram). The Transformer is the best model on \texttt{<eos>}
(\ceEosTransformer{} vs.\ \ceEosLSTM) and the worst on \texttt{<unk>} (\ceUnkTransformer{} vs.\ \ceUnkLSTM).
\emph{Cost} (Table~\ref{tab:p1c}). The 4-layer Transformer needs \trainMinTransformer{}\,min of training
against \trainMinLSTM{} for the LSTM and reaches a worse full-history perplexity with 3$\times$ the core
parameters. At inference the LSTM scores the test stream at \scoreTokSLSTM{} tokens/s against
\scoreTokSTransformer{} for the sliding-window Transformer, and samples at \genMsLSTM{} vs.\
\genMsTransformer{}\,ms per token (no KV cache). The tanh RNN trains as fast as the LSTM but scores at only
\scoreTokSRNN{} tokens/s at batch~1. The $n$-gram fits in 4\,s and scores at \scoreTokSNgram{} tokens/s, but its 156\,MB of tables
exceed any checkpoint and sampling its full distribution costs \genMsNgram{}\,ms per token.
\emph{Generated text} (Table~\ref{tab:p1gen}, Appendix~A). Under greedy decoding the RNN and the 1-layer
Transformer fall into loops (\repGreedyRNN\% and \repGreedyTransformerOne\% of 4-grams repeat an earlier
4-gram; $n$-gram \repGreedyNgram\%), while the LSTM and 4-layer Transformer repeat far less
(\repGreedyLSTM\%, \repGreedyTransformer\%). Top-$k$ sampling removes repetition for every model
($\le$\repTopkTransformerOne\%), so sampled continuations differ in local fluency and consistency instead:
the LSTM and Transformer hold a topic across the 50 tokens (a ``route description'' section after the
highway prompt, the Ming court after the 1644 prompt), the RNN drifts within a sentence and the $n$-gram
within a few tokens.

\paragraph{Explanations, separated from the observations.}
\emph{Context length, not architecture, explains LSTM $>$ Transformer.} The reversal under the
context-reset protocol is direct evidence: at $\le$128 tokens the Transformer is the better predictor,
and the LSTM's overall lead comes from information carried beyond 128 tokens (WikiText articles are
thousands of tokens long). \emph{Parameter count does not explain the ranking.} The 1-layer Transformer has about the
LSTM's core size (\coreMTransformerOne{}M vs.\ \coreMLSTM{}M) and is \deltaOneLayerVsLSTM{} points
worse, while the 4-layer model has 3$\times$ the LSTM's core and is still worse with full history; depth
matters for the Transformer, but more parameters did not buy a better model at this budget.
\emph{Training budget conditions every conclusion.} Nothing converged in 12 epochs and the ranking at
epoch~3 (RNN first) differs from that at epoch~12, so we re-ran the LSTM and Transformer for 24
epochs (last column of Table~\ref{tab:p1q}, dashed in Fig.~\ref{fig:p1curves}): the LSTM early-stopped
with its best validation at epoch \bestEpochLongLSTM{} (test PPL \pplLongLSTM), the Transformer was still
improving at epoch 24 (\pplLongTransformer), and under context reset the Transformer still leads
(\pplLongBlocksTransformer{} vs.\ \pplLongBlocksLSTM), so both conclusions survive a doubled budget. \emph{The $n$-gram's edge on rare words} is
expected from Kneser--Ney's continuation counts, which model how many contexts a rare type appears in,
whereas an embedding row updated at most ten times has barely moved from its initialisation
(a hypothesis; Part~II tests it with pre-trained rows). \emph{The RNN's slow scoring} is an MPS implementation
effect (no fused batch-1 kernel), not an architectural one, and the greedy loops of the RNN and 1-layer Transformer are consistent with their weaker modelling of the
recent context rather than with decoding, since sampling removes them.

\section{Part II: input-embedding ablation on the LSTM}
\begin{table}[t]\centering\small
\inputifexists{tables/part2_factorial.tex}
\caption{Part II: the LSTM of Part I with six input-embedding settings (rows 1, 4 and 6 are the required
comparison; test PPL $\pm$ half the range over seeds 0 and 1 where two seeds were run). Cov.: vocabulary
types with a pre-trained vector; ``val ep.\ 3/6'': validation PPL after 3 and 6 epochs (speed of convergence);
$\|g\|$: mean pre-clip gradient norm; clipped: share of steps clipped; row norm: mean $\ell_2$ norm of the
table after training.}
\label{tab:p2}
\end{table}
\begin{figure}[t]\centering
\begin{subfigure}[t]{0.49\linewidth}\centering\figifexists{part2_curves.pdf}{width=\linewidth}\caption{Validation curves (thin: seed 1).}\end{subfigure}\hfill
\begin{subfigure}[t]{0.49\linewidth}\centering\figifexists{part2_buckets.pdf}{width=\linewidth}\caption{Change in test CE vs.\ random-trainable by target frequency (hatched: frozen); the frozen random table, worse by 0.03--0.64 nats on every bucket, is omitted for scale.}\end{subfigure}
\caption{Part II learning behaviour and where each embedding setting helps or hurts.}
\label{fig:p2}
\end{figure}
\paragraph{Design.}
The LSTM of Part~I is retrained with its $300$-d input table replaced by (b) skip-gram vectors trained
by us on the WikiText-2 train stream only (gensim, negative sampling; \skipgramMinutes{}\,min on the CPU,
reported separately from LM training) and (c) public GloVe-6B vectors \citep{pennington2014glove}, both
kept entirely frozen; (a) is the Part~I model itself. Because the required three cells confound
\emph{trainability} with \emph{information content}, three cheap extra cells complete a $3\times2$
factorial: the same two pre-trained tables left trainable, and a frozen random table (no information at
all). Skip-gram covers every type (\texttt{min\_count}=1). GloVe covers \gloveCoverageTypes\% of the
types and \gloveCoverageTokens\% of train tokens; the \gloveMissing{} missing types include the most
frequent tokens of all (\texttt{\gloveMissingExamples}), and they are initialised
$\mathcal{N}(0,\sigma_d)$ with $\sigma_d$ the per-dimension standard deviation of the found rows (a
frozen zero row for \texttt{<eos>} would cripple the model). The tables differ in scale: mean row norm
\rowNormGloveTable{} (GloVe), \rowNormSkipgramTable{} (skip-gram) and about 1.0 for the random
initialisation; nothing is rescaled, which is a disclosed confound. Freezing is verified by hashing the
table before and after training. The required cells are run with seeds 0 and 1.

\paragraph{Observations.}
\emph{Prediction quality} (Table~\ref{tab:p2}). Both frozen pre-trained tables beat the jointly learned
random table: skip-gram \pplLSTMSkipFrozen{} and GloVe \pplLSTMGloveFrozen{} against \pplLSTMRand{} test
PPL, while the frozen random table collapses to \pplLSTMRandFrozen. Letting the pre-trained tables train
improves them further, to \pplLSTMSkipTrain{} and \pplLSTMGloveTrain, the two best language models of
this report. The two seeds of the random-trainable cell differ by $2\times$\pplRangeLSTMRand{} points,
which is the noise floor; every contrast discussed below is several times larger except the
skip-gram/GloVe difference among trainable tables (\cSkipVsGloveTrain{} points).
\emph{Convergence} (Fig.~\ref{fig:p2}a, columns ``val ep.\ 3/6''). The ordering is already visible after
three epochs (GloVe-trainable \valEpThreeLSTMGloveTrain{} vs.\ random \valEpThreeLSTMRand{} vs.\ random-frozen
\valEpThreeLSTMRandFrozen), and the frozen skip-gram and GloVe cells pass the random cell's final
12-epoch perplexity at epochs \epochReachLSTMSkipFrozen{} and \epochReachLSTMGloveFrozen{}; the
random-frozen cell never does. \emph{Stability.} Mean pre-clip gradient norms
(\gnMeanLSTMRand--\gnMeanLSTMRandFrozen) and the share of clipped steps ($\le$0.4\%) are
indistinguishable across cells; the only stability
differences are isolated spikes (a maximum norm of \gnMaxLSTMRandSeedOne{} in one seed of the random cell
and \gnMaxLSTMRandFrozen{} for the frozen random table, against $\le$\gnMaxLSTMGloveTrain{} elsewhere).
\emph{Generalisation.} Frozen tables fit the training stream less but generalise relatively better:
the validation--train perplexity gap at epoch~12 is \gapLSTMGloveFrozen{} (GloVe frozen) and
\gapLSTMSkipFrozen{} (skip-gram frozen) against \gapLSTMRand{} for the random table and
\gapLSTMGloveTrain{}/\gapLSTMSkipTrain{} once the pre-trained tables train, consistent with 8.7\,M fewer
trainable parameters. \emph{Where the gain sits} (Fig.~\ref{fig:p2}b). The benefit of a frozen pre-trained table is
concentrated on rare words: on targets seen at most 10 times in training, frozen GloVe and skip-gram cut
the cross-entropy by \gainRareLSTMGloveFrozen{} and \gainRareLSTMSkipFrozen{} nats relative to the random
table, the gain shrinks with frequency and vanishes on words seen more than 10k times
(\gainTopLSTMGloveFrozen{} and \gainTopLSTMSkipFrozen{} nats). Making the tables trainable shifts the
pattern: trainable GloVe gains less on the rarest words (\gainRareLSTMGloveTrain) but is the only
setting that also improves the most frequent ones (\gainTopLSTMGloveTrain), and the frozen random
table loses on every bucket.
\emph{Generated text} (Appendix~A). Sampled continuations from the pre-trained cells are locally
fluent WikiText (section headings, unit conversions such as ``2 -- 1 km ( 8 @.@ 7 mi )''), whereas the
frozen random table degenerates into \texttt{<unk> @-@ <unk>} chains (top-$k$ distinct-2
\distinctTopkLSTMRandFrozen{} vs.\ \distinctTopkLSTMRand{} for the Part~I LSTM). Under greedy decoding
repetition is \emph{not} monotone in perplexity: the best cells loop more often
(\repGreedyLSTMGloveTrain\% of 4-grams for trainable GloVe, \repGreedyLSTMSkipFrozen\% for frozen
skip-gram) than the Part~I LSTM (\repGreedyLSTMRand\%), an observation these experiments do not explain.

\paragraph{Trainability versus extra data.}
The factorial separates the two factors the brief asks about. \emph{Trainability} matters enormously
when the table carries no information (frozen vs.\ trainable random: \cRandFrozenVsTrain{} points) and
only moderately once it does (skip-gram: \cSkipFrozenVsTrain{}; GloVe: \cGloveFrozenVsTrain).
\emph{Information in the table} is worth more than trainability: replacing the trainable random table
by a frozen skip-gram table trained on the very same 2\,M tokens gains \cRandVsSkipFrozen{} points, for
\skipgramMinutes{}\,min of CPU time (under a tenth of one LM run). \emph{Extra data} (GloVe's 6\,B
tokens against our 2\,M) buys a further \cSkipVsGloveFrozen{} points frozen and \cSkipVsGloveTrain{}
trainable, i.e.\ less than having any sensible pre-trained table at all. Two caveats pull GloVe's figures in opposite directions: its most frequent missing types
(\texttt{<eos>}, \texttt{<unk>}, \texttt{@-@}) sit frozen at random vectors, a handicap, while GloVe-6B
was trained on Wikipedia 2014, which almost certainly contains the WikiText-2 source articles; and the
row-norm confound ($6\times$ longer rows) could explain part of its edge. Our reading, stated as a hypothesis: with 2\,M tokens and 12 epochs the 8.7\,M-parameter input
table is the slowest-learning part of the model (rare rows receive few updates, cf.\ the $n$-gram's
rare-word advantage in Part~I), so any procedure that gives every row a usable vector before LM training
starts is worth more than the freedom to adjust the rows during it.

\section{Part III: representations for topic classification}
\begin{table}[t]\centering\small
\inputifexists{tables/part3_main.tex}
\caption{AG News test accuracy and macro-F1 (\%) of one multinomial logistic regression on each frozen
representation, for 20,000, 2,000 and 500 labelled examples ($C$ chosen on 2,000 validation documents).}
\label{tab:p3}
\end{table}
\begin{figure}[t]\centering
\begin{subfigure}[t]{0.49\linewidth}\centering\figifexists{part3_ppl_vs_f1.pdf}{width=\linewidth}\caption{Source-LM perplexity vs.\ downstream macro-F1.}\end{subfigure}\hfill
\begin{subfigure}[t]{0.49\linewidth}\centering\figifexists{part3_learning_curve.pdf}{width=\linewidth}\caption{Macro-F1 against labelled-set size.}\end{subfigure}
\caption{Does better next-token prediction give better representations?}
\label{fig:p3}
\end{figure}
\begin{table}[t]\centering\scriptsize
\begin{minipage}[t]{0.33\linewidth}\centering\vspace{0pt}\inputifexists{tables/part3_confusion.tex}\end{minipage}\hfill
\begin{minipage}[t]{0.65\linewidth}\centering\vspace{0pt}\inputifexists{tables/part3_errors.tex}\end{minipage}
\caption{Confusion matrix (left) and the most confident errors (right) of the best LM-based representation
(\confusionRep).}
\label{tab:p3err}
\end{table}
\paragraph{Setup.}
Each AG News document (title + description, mapped to the WikiText vocabulary at \agInVocab\% token
coverage, \texttt{<eos>}-prefixed, truncated at 96 tokens) becomes one 300-d vector from a frozen
model: the mean of the top-layer states over non-pad positions, or the state at the last real token;
causal models never let padding reach a real position (unit-tested). One multinomial logistic
regression on standardised features is fitted per representation, with $C$ chosen on 2,000 held-out
training documents and the 7,600-document test split used once. Embedding-only baselines (the mean of
the raw table rows) and a TF-IDF reference over the same vocabulary ids isolate what the sequence
model adds.

\paragraph{Observations.}
\emph{Pooling.} Mean pooling beats the last state for both the LSTM (\fOneLstmRandMean{} vs.\
\fOneLstmRandLast{} macro-F1) and the Transformer (\fOneTransformerRandMean{} vs.\
\fOneTransformerRandLast): the final state of a next-token model encodes what comes next, not what the
document was about. \emph{Models.} With mean pooling the Part~I models rank 1-layer Transformer
(\fOneTransformerOnelRandMean) $>$ 4-layer Transformer (\fOneTransformerRandMean) $>$ RNN
(\fOneRnnRandMean) $>$ LSTM (\fOneLstmRandMean): the best language model is the worst feature extractor.
\emph{Embedding variants.} The LSTM with a frozen GloVe table gives \fOneLstmGloveFrozenMean{} and with
frozen skip-gram \fOneLstmSkipgramFrozenMean, far above the random table, but the raw tables alone do
better still (GloVe \fOneEmbGloveMean, skip-gram \fOneEmbSkipgramMean), and even the LSTM's own learned
input table (\fOneEmbRandLearnedMean) beats the LSTM's hidden states. \emph{Data size}
(Fig.~\ref{fig:p3}b). TF-IDF is the best representation with \nLabelsFull{} labels (\fOneTfidf) but
falls to \fOneSmallTfidf{} with \nLabelsSmall, where mean GloVe still gives \fOneSmallEmbGloveMean: dense
pre-trained features dominate the low-label regime and sparse lexical features win once labels are
plentiful; the LM-based representations sit below the embedding-only baselines at every size.
\emph{Errors} (Table~\ref{tab:p3err}, for \confusionRep). Business and Sci/Tech are confused most often
(\confBusinessSciTech{} of \nTestDocs{} test documents), then World and Business (\confWorldBusiness);
the most confident errors are label noise, sports and stock-market stories filed under \emph{World}.
Vocabulary coverage does not explain errors: the \texttt{<unk>} rate is \unkRateCorrect\% in
correctly and \unkRateWrong\% in wrongly classified documents.

\paragraph{Does the LM ranking agree with the downstream ranking?}
No. Over the \spearmanN{} LM-based mean-pooled representations the Spearman correlation between test
perplexity and macro-F1 is $\rho=\spearmanRho$ (Fig.~\ref{fig:p3}a): lower perplexity goes with
\emph{worse} classification. The disagreement has two regimes. Across architectures with the same
input table, perplexity and F1 move in opposite directions (LSTM best LM, worst features). Across
input tables with the same architecture, the three LSTM cells improve together in both metrics, but
both improvements come from the table, and the LSTM's states keep only part of what the table
provides (\fOneLstmGloveFrozenMean{} vs.\ \fOneEmbGloveMean{} for the raw GloVe mean). Our explanation,
a hypothesis consistent with all three parts: a next-token objective on Wikipedia prose rewards local
syntax and the last few words, so the state of a well-trained recurrent LM is specialised away from
topic; the Transformer's residual stream carries the input embeddings to the top layer and the
shallower Transformer carries them more faithfully (1 layer $>$ 4 layers), whereas LSTM gates
overwrite them. Domain shift (lowercased Wikipedia to news, \agInVocab\% coverage) hurts every LM-based
representation alike and cannot produce the reversal; fine-tuning might, but was outside the
frozen-representation question.

\section{Conclusions, limitations and use of AI tools}
On this corpus better next-token prediction did \emph{not} give better frozen representations, and
the input embedding table mattered more than the sequence model for both tasks. The LSTM is the best
full-history predictor only because it carries state beyond 128 tokens; at equal context the 4-layer
Transformer is better at a third more training time and half the scoring throughput. A frozen table
trained on the same 2\,M tokens in \skipgramMinutes{}\,min beats jointly learned embeddings, GloVe adds
a modest further gain, trainability a little on top, and the gains sit on rare words. Downstream,
mean-pooled raw GloVe (\fOneEmbGloveMean) beats every LM-based feature and TF-IDF wins with enough
labels (\fOneTfidf).
\emph{Limitations.} One corpus, small models and a 12-epoch budget at which nothing had converged
(the 24-epoch check in Table~\ref{tab:p1q} bounds this); MPS training is not bit-reproducible
(seed-to-seed spread about one perplexity point); GloVe's training data overlaps WikiText's source
articles and its row norms differ from the other tables; Part~III uses frozen models only, and AG
News carries visible label noise. \emph{Use of AI tools.} Claude Code (Anthropic, Claude Fable~5.1)
\citep{claudecode2026} drafted the experiment design, the code, the scripts that generate every table
and figure from the result files, and edits of this text; all experiments were run by the author and the
analysis is the author's own (see \texttt{AI\_USAGE.md}).

\bibliographystyle{unsrtnat}
\bibliography{refs}

\clearpage
\appendix
\section{All generated continuations}
\inputifexists{tables/appendix_generations.tex}
\begin{table}[h]\centering\small\inputifexists{tables/part1_repetition.tex}
\caption{Repetition statistics averaged over the five prompts.}\end{table}
\section{Embedding nearest neighbours}
\begin{table}[h]\centering\small\inputifexists{tables/appendix_neighbours.tex}
\caption{Five nearest neighbours (cosine) of probe words in the two pre-trained tables.}\end{table}
\section{Full hyper-parameters and reproduction}
\begin{table}[h]\centering\small\inputifexists{tables/appendix_hyperparams.tex}
\caption{Every setting shared by the runs; per-run differences are the model/embedding blocks in \texttt{configs/}.}\end{table}
Reproduction: \texttt{run\_all.sh} in the repository re-creates every number (about 5\,h on the M4);
\texttt{README.md} documents data access, preprocessing, dependencies, commands and seeds.
\end{document}
